\chapter{Cơ sở lý thuyết}
\label{chap:theory}

Chương này trình bày nền tảng lý thuyết của tất cả các phương pháp được nhóm triển khai, theo đúng
ba lớp tiếp cận đã giới thiệu. Với mỗi phương pháp, nhóm cố gắng tái hiện trọn vẹn mạch suy nghĩ:
\emph{động lực} dẫn tới nó, \emph{trực giác} đằng sau, \emph{công thức toán học} đúng như được cài
đặt, cùng những \emph{ưu và nhược điểm} so với các lựa chọn khác. Phần cuối chương tổng hợp tất cả
thành \emph{pipeline đề xuất} của nhóm.

\section{Khung chung và ký hiệu}
\label{sec:common}

\paragraph{Bộ trích đặc trưng (backbone).} Ở Track 1, nhóm dùng \textbf{ResNet-18} với một sửa đổi
quan trọng dành cho ảnh nhỏ: thay ``stem'' gốc của ImageNet (conv $7\times7$ stride 2 $+$ maxpool)
bằng một conv $3\times3$ stride 1 và bỏ hẳn maxpool. Lý do là stem gốc thu nhỏ ảnh $32\times32$
xuống còn $4\times4$ trước khi đặc trưng kịp hình thành, làm hỏng tín hiệu trên CIFAR. Đây là giao
thức chuẩn của CIFAR-LT (\texttt{pretrained=False}, huấn luyện từ đầu), với đầu ra là vector đặc
trưng $512$ chiều sau global average pooling. Cấu hình thay thế \texttt{pretrained=True} giữ nguyên
stem ImageNet và đòi hỏi ảnh $224$, nên chỉ được dùng cho bảng phụ.

\paragraph{Chọn mô hình theo balanced accuracy.} Một quyết định nền tảng khác là checkpoint tốt nhất
luôn được chọn theo \textbf{balanced accuracy trên tập validation} thay vì accuracy thô. Vì
validation vẫn giữ hình dạng long-tail, accuracy thô sẽ thiên về head và vô tình \emph{trừng phạt}
chính các phương pháp nhận-biết-đuôi. Quy ước này được áp dụng đồng nhất cho mọi phương pháp huấn
luyện nhằm bảo đảm so sánh công bằng.

\section{Track 1 --- Huấn luyện từ đầu: can thiệp ở tầng nào giúp đuôi?}
\label{sec:track1}

Track 1 khảo sát một câu hỏi có cấu trúc: nếu cố định backbone và chỉ can thiệp ở những \emph{tầng}
khác nhau --- mất mát, bộ phân loại, biểu diễn hay dữ liệu --- thì đuôi cải thiện ra sao?
Bảng~\ref{tab:track1-map} liệt kê các phương pháp theo tầng can thiệp.

\begin{table}[H]
\centering
\caption{Các phương pháp Track 1 phân theo tầng can thiệp.}
\label{tab:track1-map}
\begin{tabular}{lll}
\toprule
\textbf{Phương pháp} & \textbf{Tầng can thiệp} & \textbf{Ý tưởng cốt lõi} \\
\midrule
\texttt{baseline} & --- & ResNet-18 $+$ cross-entropy (mốc tham chiếu) \\
\texttt{balanced\_softmax} & mất mát & cộng log class-prior vào logit \\
\texttt{decoupling} (cRT) & bộ phân loại & học đặc trưng rồi huấn luyện lại head cân bằng \\
\texttt{supcon} & biểu diễn & học tương phản có giám sát $+$ cRT \\
\texttt{cmo} & dữ liệu & CutMix thiên-đuôi $+$ Balanced Softmax \\
\texttt{meta} & few-shot & episodic Prototypical Network \\
\bottomrule
\end{tabular}
\end{table}

\subsection{Baseline (Cross-Entropy)}
Mốc tham chiếu là ResNet-18 gắn một đầu tuyến tính, huấn luyện bằng cross-entropy chuẩn
$\mathcal{L}_{\text{CE}}=-\frac{1}{B}\sum_i \log p_{y_i}(x_i)$ với $p=\mathrm{softmax}(z)$.
Cross-entropy ngầm giả định prior lớp đồng đều; trên dữ liệu lệch, chuẩn trọng số $\lVert w_c\rVert$
tăng theo tần suất lớp, làm logit head bị thổi phồng đến mức lớp tail gần như không bao giờ thắng
được phép argmax. Đây chính là hiện tượng mà các phương pháp tiếp theo tìm cách khắc phục.

\subsection{Balanced Softmax (can thiệp ở mất mát)}
\paragraph{Động lực.} Nếu gốc rễ của vấn đề là mô hình ngầm giả định prior đồng đều, thì cách rẻ
nhất để sửa là \emph{nói cho nó biết prior thật} và bù trừ ngay trong logit. Balanced Softmax (Ren
và cộng sự, NeurIPS 2020) làm đúng điều đó.
\paragraph{Công thức.} Gọi $\pi_c = n_c/\sum_j n_j$ là prior lớp ước lượng từ tần suất train. Logit
được dịch bằng log-prior trước khi đưa vào cross-entropy:
\begin{equation}
\mathcal{L}_{\text{BS}} = \mathrm{CE}\big(z + \log \boldsymbol\pi,\; y\big),
\qquad \log\pi_c = \log\frac{n_c}{\sum_j n_j}.
\label{eq:balsoftmax}
\end{equation}
Trong cài đặt, $\log\boldsymbol\pi$ được lưu như một buffer và chỉ cộng vào lúc huấn luyện; khi test
ta dùng logit thô (giả định test cân bằng). Trực giác là lớp head có $\log\pi_c$ lớn (ít âm) nên bị
``phạt'' nhiều hơn, còn lớp tail được ``trợ giá'', nhờ vậy gradient được cân bằng lại.
\paragraph{Ưu/nhược.} Phương pháp này không thêm tham số, không đổi kiến trúc, chỉ thay một dòng
tiêu chí --- vừa rất rẻ vừa hiệu quả. Hạn chế của nó là không tác động tới phần \emph{biểu diễn}:
nếu đặc trưng của lớp tail vốn đã nghèo nàn vì quá ít dữ liệu thì việc điều chỉnh logit cũng chỉ
giúp được tới một giới hạn nhất định.

\subsection{Decoupling / cRT (can thiệp ở bộ phân loại)}
\paragraph{Động lực.} Kang và cộng sự (ICLR 2020) chỉ ra rằng trên long-tail, \emph{đặc trưng học
được thực ra đã đủ tốt}, cái lệch nằm ở \emph{bộ phân loại}. Từ nhận định đó, họ đề xuất tách
(decouple) thành hai pha: học biểu diễn trước, rồi huấn luyện lại riêng đầu phân loại.
\paragraph{Thuật toán.} Pha thứ nhất huấn luyện baseline bình thường trên dữ liệu lệch để thu được
một encoder tốt. Pha thứ hai \emph{đóng băng encoder}, thay đầu tuyến tính bằng một đầu mới khởi tạo
lại, rồi huấn luyện lại trong ít epoch (\emph{class-balanced re-training}, cRT) với \textbf{bộ lấy
mẫu cân bằng lớp} có trọng số tỉ lệ $1/n_c$. Một chi tiết cài đặt then chốt là khi cRT, encoder được
đặt ở chế độ \texttt{eval} để \emph{đóng băng thống kê BatchNorm} (cờ \texttt{eval\_encoder}); nếu
không, thống kê BN sẽ trôi theo phân bố cân bằng và làm hỏng đặc trưng đã đóng băng.
\paragraph{Ưu/nhược.} Cách làm này tách bạch rõ nguồn thiên lệch và rất rẻ ở pha hai. Tuy nhiên chất
lượng cuối cùng vẫn bị khoá bởi đặc trưng học ở pha một: nếu đuôi quá nghèo, việc cân bằng lại head
cũng không thể tạo ra đặc trưng mới.

\subsection{Supervised Contrastive (SupCon, can thiệp ở biểu diễn)}
\paragraph{Động lực.} Nếu một phần nút thắt nằm ở \emph{biểu diễn}, ta nên học một không gian đặc
trưng nơi các mẫu cùng lớp tụm lại và khác lớp tách xa, độc lập với đầu phân loại. SupCon (Khosla và
cộng sự, 2020) là một hàm mất mát tương phản có giám sát phục vụ đúng mục đích này.
\paragraph{Công thức.} Với mỗi mẫu neo $i$ trong batch (embedding $z$ đã chuẩn hoá $L_2$, nhiệt độ
$\tau$), gọi $P(i)$ là tập mẫu cùng lớp:
\begin{equation}
\mathcal{L}_{\text{SupCon}} = \sum_{i}\frac{-1}{|P(i)|}\sum_{p\in P(i)}
\log\frac{\exp(z_i^\top z_p/\tau)}{\sum_{a\neq i}\exp(z_i^\top z_a/\tau)}.
\label{eq:supcon}
\end{equation}
Cài đặt sử dụng biến thể \emph{two-crop} (mỗi ảnh hai phép tăng cường), một đầu chiếu (projection
head) dạng MLP $512\!\to\!512\!\to\!128$ và $\tau=0.07$. SupCon \emph{chỉ} pretrain trên tập train
để tránh rò rỉ validation; sau pretrain, encoder được đóng băng và áp cRT để tạo đầu phân loại cân
bằng.
\paragraph{Ưu/nhược.} SupCon có tiềm năng tạo ra biểu diễn mạnh hơn, nhưng thực nghiệm cho thấy một
hạn chế rõ: trên đuôi cực ngắn, số cặp dương của lớp tail quá ít nên SupCon khó hình thành cụm tốt
cho tail (kết quả thực nghiệm sẽ xác nhận điều này).

\subsection{CMO / Mixup / CutMix (can thiệp ở dữ liệu)}
\paragraph{Động lực.} Nếu dữ liệu tail quá thiếu, một hướng tự nhiên là \emph{tổng hợp thêm} ví dụ
tail từ những gì đang có. Mixup và CutMix là hai phép tăng cường trộn ảnh, còn CMO (Context-rich
Minority Oversampling, Park và cộng sự, CVPR 2022) là biến thể thiên-đuôi.
\paragraph{Công thức.} Với hệ số trộn $\lambda\sim\mathrm{Beta}(\alpha,\alpha)$, ba biến thể được
định nghĩa như sau:
\begin{itemize}
  \item \emph{Mixup}: $\tilde x = \lambda x_i + (1-\lambda)x_j$;
  \item \emph{CutMix}: dán một mảnh chữ nhật của $x_j$ lên $x_i$, với $\lambda$ bằng tỉ lệ diện tích
        giữ lại;
  \item \emph{CMO}: tương tự CutMix nhưng \emph{ảnh nền $x_i$ lấy từ luồng instance-balanced} còn
        \emph{mảnh dán $x_j$ lấy từ luồng class-balanced} (giàu lớp hiếm); nhờ vậy lớp hiếm được dán
        lên nhiều bối cảnh head đa dạng, tạo ra các ví dụ tail mới giàu ngữ cảnh.
\end{itemize}
Cả ba đều dùng mất mát hai-nhãn-mềm $\lambda\,\mathrm{CE}(\tilde x,y_i)+(1-\lambda)\mathrm{CE}(\tilde x,y_j)$.
Trong project, \texttt{cmo} là một phương pháp \emph{gộp}: CMO ($\alpha=1$, xác suất trộn $0.5$) kết
hợp với Balanced Softmax làm tiêu chí, nên nó can thiệp đồng thời ở cả tầng \emph{dữ liệu} và tầng
\emph{mất mát}.
\paragraph{Ưu/nhược.} CMO làm giàu đuôi một cách trực tiếp mà không cần mạng phụ, và thực nghiệm cho
thấy đây là phương pháp from-scratch mạnh nhất. Dù vậy, các ảnh trộn vẫn bị giới hạn bởi 5 ảnh gốc
của lớp tail nên không thể tạo ra ngữ nghĩa thực sự mới.

\subsection{Meta-learning --- Prototypical Network (trục few-shot)}
\paragraph{Động lực.} Nếu vấn đề cốt lõi là ``ít mẫu'', vậy hãy \emph{huấn luyện mô hình để giỏi học
từ ít mẫu}. ProtoNet huấn luyện theo từng \emph{episode} $N$-way $K$-shot, mô phỏng đúng kịch bản
few-shot.
\paragraph{Công thức.} Mỗi episode lấy $N$ lớp, mỗi lớp $K$ ảnh support cùng một số ảnh query.
Prototype của lớp $c$ là trung bình đặc trưng support, $\mu_c=\frac{1}{K}\sum_{x\in S_c}\phi(x)$, và
query được phân loại theo \emph{khoảng cách Euclid bình phương âm}
$\text{logit}_c(x)=-\lVert\phi(x)-\mu_c\rVert_2^2$ rồi đưa qua cross-entropy. Cấu hình dùng là
$5$-way $5$-shot $15$-query với Adam $\text{lr}=10^{-3}$. Khi đánh giá đủ $100$-way, prototype được
tính trên toàn bộ train rồi gán theo nearest-prototype.
\paragraph{Ưu/nhược.} Vì phân loại theo khoảng cách nên ProtoNet \emph{bất biến} với chuẩn trọng số,
qua đó tránh được kiểu thiên lệch head của đầu tuyến tính. Mặt khác, ProtoNet vốn thiết kế cho
few-way nên hoạt động yếu ở $100$-way (đúng như kỳ vọng); do đó nhóm báo cáo nó như một nhánh
\emph{bonus} trên trục few-shot chứ không kỳ vọng nó dẫn đầu bảng $100$-way.

\section{Track 2 --- Tái dùng checkpoint, không huấn luyện lại}
\label{sec:track2}

Track 2 trả lời câu hỏi: liệu có thể vắt thêm điểm số mà không cần huấn luyện lại? Tất cả các kỹ
thuật ở đây chỉ chạy ở pha inference trên những checkpoint sẵn có từ Track 1.

Kỹ thuật đầu tiên là \textbf{ensemble $+$ TTA}, lấy trung bình xác suất của nhiều mô hình
$\bar p=\frac{1}{M}\sum_m p^{(m)}$ rồi argmax. Test-time augmentation (TTA) lật ngang ảnh và trung
bình thêm; vì lỗi của các mô hình không tương quan hoàn toàn nên việc trung bình giúp giảm phương
sai. Cách này rẻ nhưng chủ yếu cải thiện head.

\textbf{Tier-fusion} trộn xác suất \emph{tham số} (đầu tuyến tính, mạnh ở head) với xác suất
\emph{prototype} (nearest-class-mean, mạnh ở tail) theo trọng số \emph{phụ thuộc nhóm-shot}:
$p_{\text{blend}}[c]=(1-w_{g(c)})\,p_{\text{param}}[c]+w_{g(c)}\,p_{\text{proto}}[c]$, với trọng số mặc
định $w_{\text{many}}{=}0,\,w_{\text{medium}}{=}0.3,\,w_{\text{few}}{=}0.8$. Đây là sự vận dụng bài
học của meta-learning --- tin prototype hơn ở đuôi --- và vì trọng số được áp theo \emph{cột-lớp}
nên ta không cần biết nhãn test.

\textbf{$\tau$-normalization} trực tiếp gỡ thiên lệch chuẩn trọng số (Kang và cộng sự, 2020) bằng
cách chuẩn hoá lại trọng số lớp $\hat w_c = w_c/\lVert w_c\rVert^{\tau}$, với $\tau$ chọn trên
validation. Giá trị $\tau{=}0$ giữ nguyên trọng số, còn $\tau{=}1$ là chuẩn hoá hoàn toàn. Vì norm
của head lớn hơn, việc chia cho norm mũ $\tau$ sẽ ``xẹp'' logit head và trả lại công bằng cho tail
mà không cần huấn luyện lại bất cứ thứ gì.

Cuối cùng, \textbf{CLIP fusion (Vision-Language)} đóng vai trò bước chuyển tiếp sang Track 3. Mô hình
thị giác nhỏ giỏi head nhưng kém tail, trong khi CLIP zero-shot nhận diện qua \emph{tên lớp} và
không hề thấy dữ liệu lệch của ta nên mạnh đều ở mọi lớp. Ta trộn lồi xác suất của hai chuyên gia,
$p=\alpha\,p_{\text{CLIP}}+(1-\alpha)\,p_{\text{vision}}$, với $\alpha$ là \emph{một số chung} chọn
trên validation theo balanced accuracy. Ở đây không có chuyện ``đoán mẫu nào hiếm''; ta chỉ trộn cho
mọi lớp rồi argmax.

\section{Track 3 --- Thích nghi mô hình nền tảng đóng băng}
\label{sec:track3}

Đây là phần trọng tâm của đề tài, xuất phát từ một luận điểm: 5 ảnh là không đủ để \emph{học từ
đầu}, nhưng đủ để \emph{thích nghi nhẹ} một mô hình nền tảng vốn đã hiểu thế giới. Mọi phương pháp
trong track này đều \textbf{đóng băng backbone} và chỉ học rất ít tham số trên \emph{đặc trưng đã
trích sẵn và cache}, nên cực nhẹ và phù hợp với môi trường Kaggle. Hình~\ref{fig:fm-pipeline} mô tả
pipeline chung.

\begin{figure}[H]
\centering
\begin{tikzpicture}[
  font=\small,
  box/.style={rectangle, rounded corners, draw, thick, align=center, minimum height=0.95cm},
  arr/.style={-{Stealth[length=2.5mm]}, thick},
]
  \node[box, fill=gray!10] (img) {Ảnh\\(train/val/test)};
  \node[box, fill=blue!10, right=1.1cm of img] (bb) {Backbone\\ \textbf{đóng băng}\\\scriptsize CLIP / DINOv2};
  \node[box, fill=blue!6, right=1.1cm of bb] (feat) {Đặc trưng\\ \textbf{cache}\\\scriptsize (trích 1 lần)};
  \node[box, fill=green!12, right=1.1cm of feat] (adapt) {Module nhẹ\\\scriptsize adapter / cache / head\\($<$1\% tham số)};
  \node[box, fill=orange!12, right=1.0cm of adapt] (out) {Logit\\ $\to$ argmax};
  \draw[arr] (img) -- (bb);
  \draw[arr] (bb) -- (feat);
  \draw[arr] (feat) -- (adapt);
  \draw[arr] (adapt) -- (out);
  \node[below=0.15cm of bb, font=\scriptsize, text=blue!50!black] {không gradient};
  \node[below=0.15cm of adapt, font=\scriptsize, text=green!40!black] {huấn luyện (Balanced Softmax)};
\end{tikzpicture}
\caption{Pipeline chung của Track 3: đóng băng backbone, trích và cache đặc trưng một lần, rồi chỉ
huấn luyện một module nhẹ ($<$1\% tham số) bằng loss nhận-biết-đuôi.}
\label{fig:fm-pipeline}
\end{figure}

\subsection{CLIP zero-shot và prototype giàu bằng LLM (nguồn: ngôn ngữ)}
\paragraph{CLIP zero-shot.} CLIP gắn ảnh và văn bản vào cùng một không gian. Với mỗi lớp, ta mã hoá
câu mẫu ``\texttt{a photo of a \{tên lớp\}}'' thành vector văn bản chuẩn hoá $t_c$; logit zero-shot
là độ tương đồng cosine có nhiệt độ $z_c(x)=s\cdot \hat\phi(x)^\top t_c$ với
$s=\exp(\text{logit\_scale})$. Vì không huấn luyện gì và không thấy dữ liệu lệch, CLIP không mang
thiên lệch head/tail. Nhóm sử dụng \textbf{ViT-B/32} (cấu hình nhẹ nhất).
\paragraph{Prototype giàu bằng LLM (CuPL).} Câu ``a photo of a \{\}'' khá mỏng về thông tin. Theo
tinh thần CuPL (Pratt và cộng sự, ICCV 2023), nhóm để một LLM nhỏ (\texttt{Qwen2.5-3B-Instruct})
sinh khoảng 8 mô tả thị giác cho mỗi lớp (về hình dáng, màu sắc, kết cấu...), mã hoá tất cả bằng
CLIP rồi \emph{trung bình} lại thành một prototype văn bản giàu hơn $\bar t_c$. Các mô tả này được
sinh một lần và cache, dùng cho cả zero-shot lẫn khởi tạo đầu LIFT.

\subsection{Tip-Adapter và Tip-Adapter-F (nguồn: truy hồi)}
\paragraph{Trực giác.} Ý tưởng ở đây là phân loại bằng \emph{truy hồi}: ảnh test giống ảnh train nào
trong ``bộ nhớ'' thì bầu theo lớp đó (k-NN mềm), rồi cộng với điểm zero-shot của CLIP. Đây chính là
Tip-Adapter (Zhang và cộng sự, ECCV 2022), một phương pháp \emph{không cần huấn luyện}.
\paragraph{Công thức.} Cache gồm \emph{khoá} $\mathbf{F}$ (đặc trưng train) và \emph{giá trị}
$\mathbf{L}$ (nhãn one-hot). Với truy vấn $\hat\phi(x)$:
\begin{equation}
\text{logit}_{\text{Tip}}(x) = \underbrace{s\,\hat\phi(x)^\top \mathbf{T}^\top}_{\text{zero-shot}}
+ \alpha\cdot \underbrace{\exp\!\big(-\beta(1-\hat\phi(x)\mathbf{F}^\top)\big)\,\mathbf{L}}_{\text{cache (k-NN mềm)}},
\label{eq:tip}
\end{equation}
trong đó $\beta$ điều khiển độ sắc của trọng số lân cận còn $\alpha$ là mức tin cache; cả hai được
\emph{grid-search trên validation} theo balanced accuracy ($\alpha\in\{1,2,5,\dots,100\}$,
$\beta\in\{1,2,3,5,7\}$). Một điểm tinh tế cho long-tail là cache mặc định sẽ lệch về head (head có
nhiều khoá hơn), nên nhóm dùng \textbf{cache cân bằng}: chia mỗi cột-lớp cho số mẫu của lớp đó, biến
tổng thành \emph{trung bình} để lớp hiếm có tiếng nói ngang với lớp phổ biến.
\paragraph{Tip-Adapter-F.} Đây là biến thể có huấn luyện nhẹ: ta coi các khoá cache như trọng số của
một lớp tuyến tính (khởi tạo bằng chính các khoá) rồi \emph{tinh chỉnh} vài epoch với Balanced
Softmax (AdamW, $\text{lr}=10^{-3}$, 20 epoch). Vì khởi tạo từ khoá, tại bước 0 nó \emph{đúng bằng}
Tip-Adapter không huấn luyện.

\subsection{LIFT --- thích nghi nhẹ với khởi tạo ngữ nghĩa (nguồn: vision-language)}
\paragraph{Động lực.} Thông điệp gốc của LIFT (Shi và cộng sự, ICML 2024) là \emph{fine-tune nặng
làm hại đuôi}. Nhóm giữ trọn tinh thần đó qua ba thành phần, minh hoạ ở Hình~\ref{fig:lift-arch}:
\begin{enumerate}
  \item \textbf{Residual Adapter có cổng (gate).} Adapter có dạng
        $h = x + g\cdot \text{Up}(\mathrm{ReLU}(\text{Down}(x)))$ với bottleneck $64$ chiều và cổng $g$
        \emph{khởi tạo bằng 0}. Tại bước 0, adapter là ánh xạ đồng nhất, nên mô hình \emph{đúng bằng
        CLIP zero-shot} và chỉ rời khỏi đó khi việc huấn luyện thực sự cải thiện được đuôi. Adapter chỉ
        chiếm dưới 1\% tham số so với backbone.
  \item \textbf{Cosine classifier khởi tạo từ feature văn bản.} Logit bằng nhiệt độ nhân với cosine
        giữa đặc trưng và trọng số lớp, cả hai cùng chuẩn hoá $L_2$. Trọng số được \emph{khởi tạo bằng
        prototype văn bản} của tên lớp, nhờ đó ngay cả lớp 5 ảnh cũng bắt đầu từ một prototype có nghĩa
        thay vì ngẫu nhiên. Việc dùng cosine còn khiến logit \emph{bất biến} với chuẩn đặc trưng, gỡ
        thêm một nguồn thiên lệch.
  \item \textbf{Loss logit-adjusted.} Mô hình được huấn luyện bằng chính \texttt{BalancedSoftmaxLoss}
        (công thức~\ref{eq:balsoftmax}) để đuôi không bị head lấn át, với AdamW, $\text{lr}=10^{-3}$,
        weight decay $10^{-2}$, $50$ epoch, và chọn epoch tốt nhất theo balanced val.
\end{enumerate}
Cần nói thêm rằng đây là biến thể LIFT ở \emph{mức đặc trưng} (adapter chạy trên feature đã cache)
nhằm giữ độ nhẹ; tuy vậy thông điệp ``đóng băng backbone $+$ học $<$1\% tham số $+$ loss
nhận-biết-đuôi'' vẫn được bảo toàn.

\begin{figure}[H]
\centering
\begin{tikzpicture}[
  font=\small,
  box/.style={rectangle, rounded corners, draw, thick, align=center, minimum height=0.85cm, minimum width=1.7cm},
  arr/.style={-{Stealth[length=2.5mm]}, thick},
]
  \node[box, fill=blue!6] (f) {feature $x$\\\scriptsize (đóng băng)};
  \node[box, fill=green!12, right=1.0cm of f] (down) {Down\\\scriptsize $D\to 64$};
  \node[box, fill=green!12, right=0.5cm of down] (relu) {ReLU};
  \node[box, fill=green!12, right=0.5cm of relu] (up) {Up\\\scriptsize $64\to D$};
  \node[circle, draw, thick, fill=yellow!20, right=0.6cm of up] (gate) {$\times g$};
  \node[circle, draw, thick, right=0.6cm of gate] (add) {$+$};
  \node[box, fill=orange!12, right=0.8cm of add] (cos) {Cosine head\\\scriptsize init $=t_c$};
  \node[box, fill=red!8, right=0.7cm of cos] (loss) {Balanced\\Softmax};

  \draw[arr] (f) -- (down);
  \draw[arr] (down) -- (relu);
  \draw[arr] (relu) -- (up);
  \draw[arr] (up) -- (gate);
  \draw[arr] (gate) -- (add);
  \draw[arr] (add) -- (cos);
  \draw[arr] (cos) -- (loss);
  \draw[arr] (f.south) to[out=-40,in=-140] (add.south);
  \node[below=0.05cm of gate, font=\scriptsize, text=orange!60!black] {$g{=}0$ lúc đầu};
\end{tikzpicture}
\caption{Kiến trúc LIFT: residual adapter có cổng ($g{=}0$ khởi tạo $\to$ bắt đầu từ zero-shot) $+$
cosine head khởi tạo từ feature văn bản, huấn luyện bằng Balanced Softmax. Backbone đóng băng.}
\label{fig:lift-arch}
\end{figure}

\subsection{DINOv2-LIFT --- mô hình nền tảng thị giác thứ hai (nguồn: thị giác)}
\paragraph{Động lực.} CLIP là một mô hình vision-language; câu hỏi đặt ra là nếu mượn một mô hình
nền tảng \emph{thị giác thuần} được huấn luyện \emph{tự giám sát} (không ngôn ngữ) thì liệu nó có
bắt được những chi tiết mà CLIP bỏ lỡ? Để trả lời, nhóm dùng \textbf{DINOv2}
(\texttt{dinov2\_vits14}, ViT-S/14, đặc trưng CLS $384$ chiều). Vì DINOv2 không có bộ mã hoá văn bản
nên cũng không có đường zero-shot; do đó đầu cosine của LIFT được khởi tạo bằng
\textbf{nearest-class-mean} (trung bình đặc trưng train mỗi lớp) thay cho prototype văn bản. Phần
còn lại --- adapter và Balanced Softmax --- giữ nguyên như LIFT, biến đây thành một thí nghiệm so
sánh trực tiếp: cùng pipeline LIFT, chỉ đổi backbone.

\subsection{Sinh đặc trưng bằng diffusion và feature mixup (nguồn: sinh dữ liệu / augment)}
\paragraph{Feature diffusion (LDMLR-style).} Thay vì sinh \emph{ảnh}, nhóm huấn luyện một DDPM nhỏ
\emph{trên không gian đặc trưng}: một MLP khử nhiễu có điều kiện theo lớp và bước thời gian (lịch
$\beta$ tuyến tính, $100$ bước), học dự đoán nhiễu theo MSE. Khi sinh, ta lấy mẫu ngược từ Gauss để
\emph{tổng hợp đặc trưng cho lớp tail} đến một mức mục tiêu (mặc định là trung vị số mẫu), rồi huấn
luyện LIFT trên hỗn hợp đặc trưng \emph{thật $+$ tổng hợp}. Ý tưởng là cân bằng lại tập huấn luyện ở
mức feature, vốn rẻ hơn nhiều so với sinh ảnh.
\paragraph{Feature mixup (CMO ở mức feature).} Ở đây nhóm chuyển ý tưởng CMO vào không gian đặc
trưng: trộn lồi một batch với các đối tác lấy từ luồng \emph{thiên-đuôi} (trọng số $\propto 1/n_c$),
$\tilde x=\lambda x_a+(1-\lambda)x_b$, $\lambda\sim\mathrm{Beta}(\alpha,\alpha)$, với mất mát
hai-nhãn-mềm; cơ chế này được kích hoạt qua \texttt{train\_lift(..., mixup\_alpha>0)}. Việc so sánh
\texttt{lift\_clip\_mixup} (trộn feature thật) với \texttt{lift\_clip\_diff} (sinh feature) chính là
để trả lời: kiểu augment nào giúp đuôi nhiều hơn?

\subsection{Generalized Logit Adjustment (GLA) --- đo công bằng}
\paragraph{Động lực.} Bản thân mô hình nền tảng cũng được huấn luyện trên dữ liệu web \emph{lệch},
nên mang sẵn một thiên lệch nhãn từ pretraining. GLA (Zhu và cộng sự, NeurIPS 2023) gỡ thiên lệch
đó, qua đó tổng quát hoá Balanced Softmax sang ``prior của foundation model''.
\paragraph{Công thức.} Ta ước lượng log-bias từ \emph{trung bình dự đoán} của mô hình trên tập đánh
giá (transductive, không dùng nhãn), $b_c=\log\big(\frac{1}{N}\sum_i \mathrm{softmax}(z(x_i))_c\big)$,
rồi điều chỉnh $\hat z = z - \gamma\, b$ với cường độ $\gamma$ chọn trên validation. Do lưới $\gamma$
luôn chứa giá trị $0$, GLA không bao giờ làm xấu đi kết quả được chọn.

\subsection{Fusion nhận-biết-đuôi (kiểm bù trừ)}
\paragraph{Động lực.} Nếu mỗi nguồn tri thức mạnh ở một vùng khác nhau của đuôi thì việc gộp chúng
đúng cách phải tốt hơn từng nguồn riêng lẻ --- và đây cũng chính là phép kiểm \emph{bù trừ}. Fusion
ở đây tổng quát hoá tier-fusion/CLIP-fusion sang $N$ chuyên gia, với \textbf{trọng số riêng cho từng
nhóm-shot} (many/medium/few) tinh chỉnh trên validation. Vì validation thiếu lớp tail, trọng số
nhóm \emph{few} được \emph{buộc theo nhóm medium} (ghi rõ trong cài đặt, không suy đoán).
\paragraph{Hai biến thể.} Biến thể thứ nhất, \texttt{fusion\_tailaware}, tìm trọng số trộn lồi theo
lưới đơn hình (simplex grid). Biến thể thứ hai, \texttt{fusion\_greedy}, dùng \emph{greedy ensemble
selection} (Caruana và cộng sự, ICML 2004; theo tinh thần greedy soup của Wortsman và cộng sự, ICML
2022): bắt đầu từ chuyên gia đơn tốt nhất mỗi nhóm và chỉ \emph{thêm} chuyên gia nếu việc đó cải
thiện balanced accuracy của nhóm. Theo cách xây dựng, greedy luôn \emph{$\ge$ chuyên gia đơn tốt
nhất} trên validation, nên không bị tụt ở many/medium như lưới có thể gặp khi một chuyên gia (DINOv2)
áp đảo.

\section{Pipeline đề xuất của nhóm}
\label{sec:proposed}

Tổng hợp lại, ``phương pháp đề xuất'' của nhóm không phải một thuật toán đơn lẻ mà là một
\emph{khung nghiên cứu có thiết kế}: một pipeline thống nhất trải qua bốn giai đoạn, trong đó mỗi
thành phần cũ được \emph{tái dùng} làm công cụ cho giai đoạn sau (Bảng~\ref{tab:reuse}).

\begin{table}[H]
\centering
\caption{Tái dùng có hệ thống các thành phần xuyên suốt bốn giai đoạn.}
\label{tab:reuse}
\begin{tabular}{p{4.2cm} p{8.6cm}}
\toprule
\textbf{Thành phần (xuất xứ)} & \textbf{Được tái dùng làm gì ở Track 3} \\
\midrule
Balanced Softmax (Track 1, loss) & Loss huấn luyện của \emph{mọi} expert LIFT / Tip-Adapter-F; được
GLA tổng quát hoá sang prior của foundation model. \\
CMO (Track 1, dữ liệu) & Chuyển vào không gian feature thành \texttt{lift\_clip\_mixup}. \\
Tier-fusion / CLIP-fusion (Track 2) & Tổng quát hoá thành fusion $N$-expert nhận-biết-đuôi (kiểm bù trừ). \\
Mô hình \texttt{cmo} from-scratch (Track 1) & Mốc \emph{đối chứng} ``không tri thức ngoài'' cho Phase 3. \\
\bottomrule
\end{tabular}
\end{table}

Một số thiết kế then chốt cùng lý do của chúng có thể nêu ra như sau. Trước hết là việc
\emph{đóng băng backbone và chỉ học dưới 1\% tham số}: với 5 ảnh/lớp, fine-tune nặng sẽ overfit và
``quên'' tri thức zero-shot ở đuôi --- điều sẽ được kiểm chứng bằng Ablation B --- trong khi việc
cache feature một lần khiến toàn bộ Track 3 chạy được trong vài phút trên GPU Kaggle. Tiếp theo là
\emph{khởi tạo có ngữ nghĩa}: cosine head được khởi tạo từ feature văn bản (CLIP) hoặc class-mean
(DINOv2), cho lớp 5 ảnh một điểm xuất phát hợp lý thay vì ngẫu nhiên. Bên cạnh đó, \emph{loss
nhận-biết-đuôi} (Balanced Softmax / GLA) được dùng xuyên suốt nhằm bảo đảm gradient không bị head
chi phối ở bất kỳ expert nào. Cuối cùng, \emph{mọi siêu tham số đều được chọn trên validation} chứ
không đụng tới test, để kết quả thu được là đáng tin cậy.
